\documentclass[12pt,a4paper]{article}

% --- Encoding & Language ---
\usepackage[utf8]{inputenc}
\usepackage[T5]{fontenc}
\usepackage[vietnamese]{babel}

% --- Layout ---
\usepackage[margin=2.5cm]{geometry}
\usepackage{parskip}
\setlength{\parindent}{0pt}

% --- Math ---
\usepackage{amsmath, amssymb}

% --- Tables ---
\usepackage{booktabs}
\usepackage{array}
\usepackage{tabularx}
\usepackage{xcolor}

% --- Code blocks ---
\usepackage{listings}
\lstset{
  basicstyle=\ttfamily\small,
  backgroundcolor=\color{gray!10},
  frame=single,
  framerule=0.4pt,
  rulecolor=\color{gray!50},
  breaklines=true,
  showstringspaces=false,
}

% --- Boxes ---
\usepackage{tcolorbox}
\tcbuselibrary{skins, breakable}

% --- Hyperlinks ---
\usepackage{hyperref}
\hypersetup{
  colorlinks=true,
  linkcolor=teal,
  urlcolor=teal,
}

% --- Colors ---
\definecolor{teal}{RGB}{13, 148, 136}
\definecolor{violet}{RGB}{124, 58, 237}
\definecolor{amber}{RGB}{245, 158, 11}
\definecolor{navy}{RGB}{15, 23, 42}
\definecolor{boxbg}{RGB}{240, 253, 250}

% --- Section styling ---
\usepackage{titlesec}
\titleformat{\section}{\large\bfseries\color{navy}}{}{0em}{}[\titlerule]
\titleformat{\subsection}{\normalsize\bfseries\color{teal}}{}{0em}{}

% --- Header / Footer ---
\usepackage{fancyhdr}
\pagestyle{fancy}
\fancyhf{}
\rhead{\textcolor{gray}{\small Lý Thuyết Đồ Thị}}
\lhead{\textcolor{gray}{\small Nhóm: Codemax-cyber · bk-nguyen · levanh1711-cyber}}
\cfoot{\thepage}

% ============================================================
\begin{document}

% --- Title ---
\begin{center}
  {\LARGE\bfseries Lý Thuyết Đồ Thị}\\[0.4em]
  {\large\color{gray} Nền tảng lý thuyết đồ thị --- từ định nghĩa cơ bản đến các cấu trúc biểu diễn quan trọng}\\[0.6em]
  {\small\color{gray} Nhóm: Codemax-cyber \textbullet{} bk-nguyen \textbullet{} levanh1711-cyber}\\
  {\small\color{gray} \url{https://github.com/Codemax-cyber/Graph}}
\end{center}

\vspace{0.5em}
\hrule
\vspace{1.5em}

% ============================================================
\section{1. Đồ thị là gì?}

\textbf{Đồ thị (Graph)} $G = (V, E)$ là một cấu trúc toán học gồm:

\begin{itemize}
  \item \textbf{V (Vertices)}: Tập hợp các \textbf{đỉnh} (nodes). Ví dụ: thành phố, người dùng, trang web.
  \item \textbf{E (Edges)}: Tập hợp các \textbf{cạnh} (edges) nối các đỉnh. Ví dụ: đường đi, mối quan hệ, liên kết.
\end{itemize}

\begin{tcolorbox}[colback=boxbg, colframe=teal, fontupper=\ttfamily\small]
$G = (V, E)$ \quad $V = \{A, B, C, D\}$ \quad $E = \{(A,B),\, (B,C),\, (C,D),\, (A,D)\}$
\end{tcolorbox}

% ============================================================
\section{2. Các khái niệm cơ bản}

\begin{tabularx}{\textwidth}{>{\ttfamily\bfseries\color{teal}}p{4.5cm} X}
  \toprule
  \normalfont\textbf{Khái niệm} & \textbf{Định nghĩa} \\
  \midrule
  Đỉnh (Vertex)       & Một nút trong đồ thị, thường biểu diễn bằng một ký tự hoặc số. \\
  Cạnh (Edge)         & Kết nối giữa hai đỉnh, có thể có trọng số (weighted) hoặc không. \\
  Bậc (Degree)        & Số cạnh liên kết với một đỉnh. Ký hiệu: $\deg(v)$. \\
  Đường đi (Path)     & Dãy các đỉnh liên tiếp nhau qua các cạnh, không lặp đỉnh. \\
  Chu trình (Cycle)   & Đường đi bắt đầu và kết thúc tại cùng một đỉnh. \\
  Đồ thị liên thông   & Mọi cặp đỉnh đều có đường đi nối với nhau. \\
  Cây (Tree)          & Đồ thị vô hướng liên thông và không có chu trình. \\
  DAG                 & Directed Acyclic Graph --- đồ thị có hướng không có chu trình. \\
  Đỉnh kề (Neighbor)  & Hai đỉnh $u$, $v$ kề nhau nếu tồn tại cạnh $(u, v) \in E$. \\
  \bottomrule
\end{tabularx}

% ============================================================
\section{3. Phân loại đồ thị}

\begin{tabularx}{\textwidth}{>{\bfseries}p{5.5cm} X}
  \toprule
  \textbf{Loại đồ thị} & \textbf{Đặc điểm} \\
  \midrule
  Đồ thị vô hướng
    & Cạnh không có chiều. Nếu có cạnh $(u,v)$ thì cũng có $(v,u)$.
      Ví dụ: mạng xã hội đơn giản. \\[4pt]
  Đồ thị có hướng (Digraph)
    & Cạnh có chiều xác định. $(u,v) \neq (v,u)$.
      Ví dụ: Twitter follow, web link. \\[4pt]
  Đồ thị có trọng số
    & Mỗi cạnh gắn với một giá trị (trọng số).
      Ví dụ: khoảng cách, chi phí, thời gian. \\[4pt]
  Đồ thị không trọng số
    & Tất cả cạnh có giá trị bằng nhau (mặc định $= 1$).
      Chỉ quan tâm sự liên kết. \\[4pt]
  Đồ thị thưa (Sparse)
    & Số cạnh $|E| \ll |V|^2$. Dùng danh sách kề để biểu diễn hiệu quả. \\[4pt]
  Đồ thị dày (Dense)
    & Số cạnh $|E| \approx |V|^2$. Dùng ma trận kề để biểu diễn và truy cập nhanh. \\
  \bottomrule
\end{tabularx}

% ============================================================
\section{4. Biểu diễn đồ thị}

\subsection{4.1 Ma trận kề (Adjacency Matrix)}

Ma trận vuông $n \times n$, trong đó $A[i][j] = 1$ nếu có cạnh $(i,j)$, $= 0$ nếu không.

\begin{lstlisting}
    A  B  C  D
A [ 0  1  1  0 ]
B [ 1  0  0  1 ]
C [ 1  0  0  1 ]
D [ 0  1  1  0 ]
\end{lstlisting}

\begin{itemize}
  \item[\textcolor{teal}{+}] Kiểm tra cạnh: $O(1)$
  \item[\textcolor{red}{--}] Bộ nhớ: $O(V^2)$
\end{itemize}

\subsection{4.2 Danh sách kề (Adjacency List)}

Mỗi đỉnh lưu danh sách các đỉnh kề. Phù hợp cho đồ thị thưa.

\begin{lstlisting}
A: [B, C]
B: [A, D]
C: [A, D]
D: [B, C]
\end{lstlisting}

\begin{itemize}
  \item[\textcolor{teal}{+}] Bộ nhớ: $O(V + E)$
  \item[\textcolor{red}{--}] Kiểm tra cạnh: $O(\deg(u))$
\end{itemize}

% ============================================================
\section{5. Độ phức tạp so sánh}

\begin{tabularx}{\textwidth}{X >{\centering\arraybackslash\ttfamily\color{amber}}p{3.5cm} >{\centering\arraybackslash\ttfamily\color{teal}}p{3.5cm}}
  \toprule
  \textbf{Thao tác} & \textbf{Ma trận kề} & \textbf{Danh sách kề} \\
  \midrule
  Kiểm tra cạnh $(u,v)$   & $O(1)$       & $O(\deg(u))$ \\
  Tìm tất cả đỉnh kề $u$  & $O(V)$       & $O(\deg(u))$ \\
  Thêm cạnh               & $O(1)$       & $O(1)$ \\
  Xóa cạnh                & $O(1)$       & $O(\deg(u))$ \\
  Bộ nhớ                  & $O(V^2)$     & $O(V+E)$ \\
  \bottomrule
\end{tabularx}

% ============================================================
\vspace{2em}
\hrule
\vspace{0.8em}
{\small\color{gray}
Tài liệu được tạo tự động từ dự án web \textit{Lý Thuyết Đồ Thị}.\\
Xem trực quan hóa tương tác tại: \\url{https://github.com/Codemax-cyber/Graph}
}

\end{document}
